\chapter{Protocolo experimental: de la señal a la cartera}
\label{chap:diseno-experimental}

Una arquitectura como la del capítulo anterior no produce un sistema hasta que alguien fija el valor
de cada una de sus variables. Este capítulo explica cómo se eligieron, comparando Rank-IC y sin que
la rentabilidad observada intervenga en ninguna decisión predictiva, y cómo la ordenación resultante
se convierte después en posiciones concretas. Las cifras que aparecen aquí documentan decisiones
---qué alternativa se comparó con cuál y qué regla resolvió el empate---; si la señal resultante es
real y si se traduce en rentabilidad son preguntas distintas, que responden los dos capítulos
siguientes.

\section{Por qué el protocolo se aplicó tres veces}
\label{sec:cadena-studies}

El protocolo no se ejecutó una vez sino tres, encadenadas, y todas las cifras que siguen son las de
la tercera pasada.

El motivo está en cómo funciona la selección secuencial. Las fases se recorren en orden y cada una
parte de lo que decidió la anterior, así que el resultado depende del punto de partida: la primera
variable se elige con todas las demás en su valor por defecto, y ya no se revisa aunque las
decisiones posteriores cambien el contexto en que aquella elección tenía sentido.

La primera pasada lo dejó claro. Su ganador se apartaba en \textbf{nueve variables} del punto de
partida del catálogo, de modo que casi todas sus decisiones se habían tomado sobre una
configuración que el propio estudio acabó descartando.

La respuesta fue encadenar los estudios: el ganador de cada pasada se convierte en el punto de
partida de la siguiente, que vuelve a recorrer las fases desde ahí. Cada pasada sólo puede mejorar o
confirmar, porque el valor vigente se conserva cuando no hay evidencia suficiente para desplazarlo.

El procedimiento se fue quedando sin cambios que hacer: nueve variables modificadas en la primera
pasada, dos en la segunda y dos en la tercera. Eso es lo que debe ocurrir si converge en lugar de
explorar al azar, y es la razón de detenerse en tres.

Encadenar tiene un coste. Multiplica el número de configuraciones evaluadas sobre la misma ventana,
y con ello la probabilidad de que la mejor cifra observada deba parte de su tamaño a la propia
búsqueda. La cadena no evita el sobreajuste: lo sistematiza y lo hace visible. Lo que protege al
resultado es la era reservada, que no participó en ninguna de las tres pasadas.

\section{Optimización secuencial y regla de decisión}

El catálogo cerrado se recorre en cuatro fases: temporal, representación, modelo y meta. Se parte
de un \textit{baseline}; en cada fase se modifica una variable y se comparan sus candidatos con el
incumbente acumulado. Así, el coste científico es una suma de alternativas, no una búsqueda
cartesiana opaca. La pasada de referencia ejecutó 46 configuraciones, y cada una queda registrada
junto con la elección que motivó.

\subsection*{Cuándo una alternativa sustituye a la que ya estaba}
\label{sec:regla-decision}

En cada paso hay una configuración \emph{en el cargo}, el incumbente, y una alternativa que aspira a
desplazarla. El procedimiento está deliberadamente sesgado hacia no cambiar: si la evidencia no
distingue con holgura entre ambas, se queda la que ya estaba. Cambiar exige demostrar algo; no
cambiar, no.

Lo primero que conviene entender es \emph{qué} se compara, porque no son dos promedios. Se compara
la misma fecha contra la misma fecha: en cada mes se mide cuánto ordena mejor o peor el candidato
que el incumbente, y lo que se juzga es esa serie de diferencias. Así una alternativa no puede ganar
por haber acertado un año excepcional que el incumbente no vivió.

Sobre esa serie, un candidato puede ganar de dos maneras.

La primera es \textbf{ganar claramente}: su diferencia media es favorable \emph{y} además gana en
más de la mitad de los meses. Hacen falta las dos condiciones, porque una media favorable construida
sobre menos de la mitad de las fechas no demuestra superioridad: significa que unos pocos meses muy
buenos compensaron a muchos malos.

La segunda es \textbf{ganar por no ser peor}. Cuando la alternativa es más simple o más
conservadora, basta con demostrar que su desventaja plausible es despreciable, es decir, que ni en
el peor escenario razonable pierde apenas nada. Prueba que no empeora, no que mejore.

Hay además dos filtros previos que un candidato debe pasar antes de llegar a esa comparación: tener
observaciones suficientes y no hundirse en ninguna era disponible, de modo que una configuración que
fracasa por completo en un régimen queda excluida aunque su promedio sea favorable.

Y cuando la diferencia es minúscula ---por debajo de 0,002--- se declara empate técnico y decide una
preferencia fijada de antemano hacia la opción más simple o más prudente. Eso no es una medición, y
por eso los valores que salen de un empate no se presentan en este trabajo como hallazgos empíricos.

Un caso real del propio estudio ilustra por qué la regla mira la serie y no el promedio. La regla de
ponderación de variables tenía una alternativa que medía \emph{mejor} ---0,1065 frente a 0,1063--- y
cuya ventaja pareada era incluso positiva, y aun así no ganó: era superior en el 49,6\,\% de las
cohortes, es decir, en menos de la mitad. No superó la puerta de dominancia y el incumbente se
conservó.

Esa separación es la garantía que sostiene todo lo demás: sólo las cuatro fases predictivas pueden
mover al ganador, y todo lo económico ocurre después de congelar la señal.

La Tabla~\ref{tab:decisiones} hace auditable esa traza. Lo que tabula no es el Rank-IC acumulado
tras cada paso ---que apenas se mueve, por lo que se explica más abajo--- sino \textbf{la
alternativa que se descartó}: cuánto Rank-IC habría costado adoptarla, ordenado de mayor a menor
coste.

\begin{table}[H]
  \centering
  \small
  \caption{Qué alternativa se rechazó en cada decisión y cuánto habría costado adoptarla.}
  \label{tab:decisiones}
  \resizebox{\textwidth}{!}{\input{tables/t05_decisiones.tex}}
  \caption*{El coste es el Rank-IC de la mejor alternativa menos el del ganador; un coste
  \emph{positivo} indica que la alternativa medía mejor y que el ganador se decidió por otra vía, lo
  que ocurre en dos filas y se comenta más abajo. No debe confundirse con la \emph{ventaja pareada},
  que promedia la diferencia cohorte a cohorte: en el método de combinación del meta-agente el coste
  es $-0{,}0071$ y la ventaja pareada $+0{,}00699$.}
\end{table}

Leída así, la tabla separa con claridad tres grupos. Cinco decisiones eran caras de equivocar:
cambiar el objetivo de regresión de rango a \textit{ranking} habría costado $-0{,}0768$ de Rank-IC,
neutralizar por sector $-0{,}0514$, reducir el conjunto de variables a \texttt{core} $-0{,}0443$,
acortar el horizonte a seis meses $-0{,}0379$ y aplicar pesos de recencia $-0{,}0332$. Son las
decisiones estructurales, y las cinco se resolvieron con holgura.

Un segundo grupo, el más numeroso ---diez de las dieciocho filas---, tiene costes de $0{,}005$ o
menos: los cuatro hiperparámetros del modelo, el retardo de observación, la winsorización, el máximo
de variables por agente y la regla que las pondera. Ahí la evidencia apenas distingue entre
alternativas, y el resultado merece decirse sin rodeos: \textbf{el sistema no debe su capacidad
predictiva al ajuste fino del modelo}.

Entre ambos extremos quedan tres decisiones intermedias ---la cadencia de observación, la ventana de
entrenamiento y la variante del meta-agente---, que también confirmaron el valor que ya traían.

El tercer grupo son las \textbf{cuatro} decisiones que la última columna marca como empate técnico:
el retardo de observación, la winsorización, la variable de régimen de mercado y la tasa de
aprendizaje. Ninguna de las cuatro puede presentarse como hallazgo empírico, y por eso figuran con
su regla al lado.

El dato que resume la pasada es que \textbf{dieciséis de las dieciocho decisiones confirmaron su
valor de partida}, y las dos que lo cambiaron lo hicieron por márgenes mínimos: la variable de
régimen de mercado por no inferioridad, sin demostrar ser mejor, y la tasa de aprendizaje con una
ventaja media de tres diezmilésimas. Esa planitud es lo esperable de una cadena que ha convergido
---la pasada parte del ganador de la anterior, así que casi todas las variables llegan ya en su
mejor valor conocido--- y es también una defensa parcial frente a la objeción de sobreajuste: un
protocolo que aceptase cualquier mejora aparente habría movido al incumbente muchas más veces,
persiguiendo ruido.

Dos filas merecen lectura aparte porque en ellas \textbf{el ganador no es el valor que mejor midió}.
En la winsorización, la alternativa medía cinco diezmilésimas por encima, la diferencia cayó dentro
del empate técnico y se impuso la preferencia por la opción más simple, que aquí es no winsorizar.
La otra es la regla de ponderación de variables, ya comentada al presentar la regla de decisión: la
alternativa medía mejor en promedio pero ganaba en menos de la mitad de las cohortes.

El catálogo cerrado es lo que hace verificable todo esto. El Anexo~\ref{anexo:catalogo} reproduce
las treinta y tres variables con su rejilla completa y el valor que tomó el ganador, de modo que
puede comprobarse que cada valor final pertenecía a su rejilla declarada y que ninguna variable
ajena aparece en la configuración. Doce de esas treinta y tres son de cartera: se ejecutan después
de congelar el ganador y no pueden modificarlo.

\section{Por qué la selección predictiva es secuencial}
\label{sec:seleccion-secuencial}

Una búsqueda cartesiana de todos los valores posibles permitiría encontrar combinaciones muy
ajustadas a la historia, dificultaría explicar de dónde surge el ganador y multiplicaría las
pruebas. El protocolo fija un orden: temporalidad, representación, modelo y meta. Cada fase hereda
el incumbent anterior, de modo que una decisión posterior no puede reabrir silenciosamente una
elección temporal. Si dos alternativas difieren menos de 0,002 en ventaja pareada, se prefiere la
más simple; si tampoco hay evidencia suficiente, se conserva el incumbent.

Ese argumento vale para el plano \textbf{predictivo} por dos motivos. El primero es estadístico: la
selección predictiva se decide sobre el Rank-IC, la misma magnitud que después se reporta como
resultado principal, de modo que cada comparación adicional infla precisamente la cifra que sostiene
el trabajo. El segundo es expositivo: una rejilla de miles de puntos devuelve un ganador que nadie
puede explicar.

Ninguno de los dos se aplica al plano de la \textbf{cartera}, y de ahí que allí el trabajo haga lo
contrario. Las reglas de cartera operan sobre una señal ya congelada, así que optimizarlas no puede
inflar la métrica predictiva; y su elección no es una escalera de decisiones, sino un problema con
interacciones fuertes que el apartado del Portfolio Study desarrolla.

Un ejemplo de cada puerta cierra el mecanismo. La cadencia mensual desplazó a la trimestral por
\emph{dominancia}: ventaja media de $+0{,}0136$ sobre 40 fechas emparejadas y mejor en el 60\,\% de
ellas, aunque el límite inferior de su intervalo rozaba cero. Un procedimiento que hubiera exigido
además un intervalo estrictamente positivo habría conservado la cadencia trimestral, y con ella
habría perdido las 117 cohortes mensuales sobre las que descansa el resto del trabajo. En sentido
contrario, el horizonte de seis meses fue rechazado con holgura: ventaja pareada de $-0{,}0375$ y un
intervalo íntegramente negativo.

\section{De rango a alfa esperado}

A partir de aquí el problema cambia de naturaleza: ya no se trata de ordenar mejor, sino de
convertir una ordenación en posiciones sin destruir por el camino la ventaja que contiene. El primer
paso es que una cartera necesita una magnitud económica, no sólo un puesto en una lista: saber que
una acción es la tercera mejor no dice si merece la pena pagar por comprarla. Sobre cohortes ya
cerradas se estima por eso una relación entre el tramo del ranking y el retorno excedente anual,
usando veinte grupos para que cada tramo conserve suficientes acciones.

\subsection*{Calibración causal y estados de respaldo}

En cada fecha de inversión el calibrador sólo mira cohortes cuyo retorno a doce meses ya terminó.
Para cada tramo del ranking calcula el exceso medio que se observó, y con esos veinte puntos ajusta
una función creciente que traduce posición en alfa esperado:
\[
 \widehat{\alpha}_{i,t}=g_t\!\left(\operatorname{rank}_{i,t}\right).
\]
Que la función sea creciente no obliga al mercado a pagar ninguna cantidad concreta. Impone sólo una
coherencia mínima: que una acción mejor ordenada no reciba menos alfa esperado que otra peor
ordenada. La magnitud sale de la historia cerrada y se expresa en puntos básicos anuales, de modo
que los umbrales operativos significan lo mismo aunque cambie el horizonte.

La estimación más local usa cohortes recientes del mismo horizonte; si no hay historia suficiente se
amplía a dieciséis trimestres y, si tampoco basta, a todo el histórico cerrado. Sólo cuando ninguna
de las tres es viable se activa una salvaguarda fija. Cada observación registra cuál de los cuatro
estados empleó, lo que evita mezclar silenciosamente un supuesto con una estimación.

\begin{table}[H]
  \centering
  \small
  \caption{Frecuencia de los estados de calibración por ventana.}
  \label{tab:calibracion-estados}
  \begin{tabular}{llrr}
\toprule
Ventana & Fuente de la curva & Snapshots & Peso en la ventana \\
\midrule
Selección 2015--2024 & Sin calibración & 15 & 12,8\% \\
Selección 2015--2024 & Salvaguarda & 6 & 5,1\% \\
Selección 2015--2024 & Era & 24 & 20,5\% \\
Selección 2015--2024 & Horizonte & 72 & 61,5\% \\
Reserva 2025--2026 & Sin calibración & 0 & 0,0\% \\
Reserva 2025--2026 & Salvaguarda & 0 & 0,0\% \\
Reserva 2025--2026 & Era & 0 & 0,0\% \\
Reserva 2025--2026 & Horizonte & 18 & 100,0\% \\
\bottomrule
\end{tabular}

\end{table}

En selección, 72 de las 117 observaciones usan la estimación del mismo horizonte, 24 amplían a la
era, 15 quedan sin calibración empírica suficiente y seis activan la salvaguarda de arranque; en la
reserva, las 18 aplican la estimación de horizonte obtenida con historia anterior. El recuento hace
visible cuándo el sistema aprendió una escala económica y cuándo operaba con una regla de seguridad.

La calibración no crea información predictiva: dos acciones conservan el orden que les dio el
meta-agente, y la curva sólo cuantifica la distancia económica necesaria para compararlas con
costes, efectivo y una posición existente. Por eso sus parámetros no participan en la selección por
Rank-IC.

\section{Construcción de cartera}

Seis variables gobiernan cómo se traduce el ranking en posiciones: cuántas acciones se mantienen a
la vez, cómo se reparte el peso entre ellas, cuánto efectivo se tolera, cuánto tiempo debe retenerse
una posición antes de poder venderla, por debajo de qué percentil se vende una que ha decaído y
cuánta deriva se consiente en los pesos antes de corregirla. El catálogo declara la rejilla de cada
una; qué valores toman se decide con el procedimiento de la Sección~\ref{sec:portfolio-study}.

Las restantes condiciones de simulación sí son constantes del entorno y no se optimizan: comisión de
10 pb, \textit{slippage} de 20 pb y SPY como benchmark, nunca como posición. Los umbrales se expresan
en puntos básicos anuales de alfa esperado para poder compararlos entre horizontes. Una rotación
requiere que la diferencia esperada cubra el coste de ida y vuelta y un margen de 50 pb; una compra
nueva usa histéresis. Estas reglas pretenden preservar señal y controlar rotación, no volver a
optimizar el modelo por rentabilidad.

La era 2025--2026 fue aislada antes de aplicar la selección. Sus observaciones no participan ni en
los pesos del meta ni en la elección de parámetros; su papel es exclusivamente de confirmación
posterior.

\subsection*{Precedencia de las reglas operativas}

Las seis variables no se aplican como filtros independientes: en cada rebalanceo hay una precedencia
explícita que impide contradicciones. Primero se actualizan señales y restricciones de permanencia;
después se evalúan las ventas económicamente elegibles; a continuación se comparan destinos para el
capital liberado; y por último se ajustan tamaños, sólo si la deriva justifica pagar el coste.

Esa jerarquía aclara una diferencia importante: la permanencia mínima puede impedir cualquier venta,
mientras que el suelo de cobertura o un alfa negativo sólo convierten una posición en candidata a
venderse. Incluso entonces debe existir un destino mejor después de costes. El índice no puede ser
ese destino, porque sólo actúa como referencia; el efectivo sí, pero dentro de su tope.

\section{Ejemplo de decisión y coste}

\begin{ejemplo}[Cuándo se rota y cuándo no]
Con horizonte de doce meses, comisión de 10 pb y \textit{slippage} de 20 pb, una rotación completa
--- vender una posición y comprar otra --- cuesta 60 pb. Si la peor posición tuviese alfa esperado
60 pb, una candidata externa 250 pb y el margen de rotación fuese 50 pb, la mejora de 190 pb
superaría los 110 pb exigidos entre coste y margen, y se rotaría. Si la candidata ofreciese 150 pb,
la mejora de 90 pb no cubriría esos 110 pb y se conservaría la posición.
\end{ejemplo}

El ranking ordena preferencias; operar exige además una ventaja neta. Las cifras del ejemplo son
inventadas ---los costes sí son los de la simulación, pero los alfas no proceden de ninguna
medición--- y la regla que ilustran es la que gobierna toda la rotación de la cartera. El efectivo,
por su parte, no se remunera, y una plaza sólo puede quedarse en caja dentro del tope declarado: si
ese tope es cero, debe reinvertirse necesariamente.

\section{El Portfolio Study: optimizar la cartera sin tocar la señal}
\label{sec:portfolio-study}

Elegir esas seis variables por su rentabilidad conocida parecería contradecir la regla que gobierna
el protocolo, pero no lo hace. Lo que la regla protege es la \emph{métrica predictiva}: si se elige
una configuración porque su Rank-IC es alto y luego se reporta ese mismo Rank-IC, la cifra está
inflada por la búsqueda. Las reglas de cartera no participan de ese circuito, porque operan sobre
una señal congelada y no pueden modificar el Rank-IC. Lo que sí hay que proteger es la era
reservada, y eso se consigue por construcción.

De ahí el \emph{Portfolio Study}. Toma el ganador del último Model Study, no reentrena nada y
recorre el producto cartesiano de las seis variables midiendo cada combinación sobre una serie
\textbf{recortada en 2024}. La distinción es fina y decisiva: no se trata de calcular el resultado
de cada cartera en la era reservada y luego ignorarlo ---si existiera, bastaría con mirarlo--- sino
de que ese resultado \emph{no llega a calcularse}. Sólo la ganadora, ya elegida, se reevalúa después
sobre la serie completa, como comprobación y nunca como criterio.

\subsection*{Qué es el Information Ratio y por qué se optimiza éste}

El criterio de selección deja de ser el exceso geométrico y pasa a ser el \emph{Information Ratio}:
el exceso medio sobre el índice dividido por la volatilidad de ese exceso. Premia la
\textbf{consistencia}, no el alfa bruto. Una cartera que bate al índice un 2\,\% todos los años
tiene un IR altísimo; otra que alterna $+15$\,\% y $-11$\,\% puede tener el mismo exceso medio y un
IR pésimo. Para un trabajo cuya pregunta es si una señal se sostiene fuera de muestra, la segunda
cartera no es equivalente a la primera aunque el promedio coincida, y optimizar por exceso bruto no
sabría distinguirlas.

Cada una de las seis mueve ese cociente por una vía distinta, y eso explica la forma de los
resultados:

\begin{itemize}
  \item El \textbf{número de posiciones} gobierna el compromiso central. Pocas concentran la señal
  pero hacen el exceso volátil; muchas lo estabilizan a costa de parecerse al índice, y una cartera
  que se parece al índice tiene poco exceso que mostrar.
  \item El \textbf{tope de efectivo} amortigua numerador y denominador a la vez: la caja no se
  remunera, así que sostenerla resta rentabilidad esperada, pero también resta varianza.
  \item El \textbf{reparto de pesos} decide si el peso sigue al alfa esperado o se distribuye por
  igual. Concentrar en las convicciones altas sube el exceso y la varianza a la vez.
  \item La \textbf{permanencia mínima} reduce fricciones, pero no de forma monótona: retener de más
  impide corregir cuando la señal cambia de opinión, y ese coste de oportunidad acaba superando al
  ahorro.
  \item El \textbf{suelo de cobertura} corta la cola inferior del ranking, forzando la venta de lo
  que cae por debajo del percentil declarado.
  \item La \textbf{tolerancia a la deriva} es pura fricción: no cambia qué se tiene, sino cada
  cuánto se corrige el peso de lo que ya se tiene.
\end{itemize}

\subsection*{Por qué cartesiano y no secuencial}

Aquí se hace lo contrario que en el plano predictivo, y no por comodidad: las seis variables
\textbf{interactúan}, de modo que recorrerlas de una en una daría una respuesta que depende del
orden en que se pregunte.

El suelo de cobertura lo ilustra bien, porque hace cosas \emph{opuestas} según cuánto efectivo se
tolere. Con tope cero, la plaza que deja una posición vendida se recompra necesariamente, así que el
suelo funciona como un mecanismo de sustitución. Con tope positivo, esa misma plaza puede quedarse
en caja, y entonces el mismo valor pasa a reducir la exposición.

Un ascenso por coordenadas sobre variables así converge a un óptimo local que depende de por dónde
se empiece; el cartesiano no. Tiene otro problema, la multiplicidad, que declara el
Capítulo~\ref{chap:limitaciones}.

\subsection{Dos búsquedas y dos multiplicidades}
\label{sec:dos-multiplicidades}

El proyecto contiene dos recuentos que no deben sumarse ni intercambiarse. La pasada predictiva de
referencia evaluó 46 configuraciones del Model Study. Esas comparaciones podían afectar qué señal
se congelaba y, por tanto, forman el universo de intentos que penaliza la inferencia predictiva y el
Deflated Sharpe asociado al resultado económico de esa señal.

El Portfolio Study evaluó 1.440 combinaciones de seis reglas sobre la selección hasta 2024. Es una
multiplicidad distinta: la señal, sus pesos y su Rank-IC permanecían idénticos en las 1.440 celdas.
Este recuento obliga a tratar el Information Ratio ganador como seleccionado y a reservar
2025--2026 para una única reevaluación, pero no aumenta retrospectivamente el número de modelos
predictivos. Mezclar ambos números produciría dos errores: atribuir a la cartera capacidad de elegir
el modelo y presentar el máximo de la rejilla como una estimación no condicionada por búsqueda.

La regla de lectura queda así fijada una sola vez. El Capítulo~\ref{chap:resultados} juzga la señal
con la multiplicidad de Model Study; el Capítulo~\ref{chap:resultados-economicos} interpreta la
cartera seleccionada frente a su distribución de 1.440 alternativas y luego mira una sola reserva.
Las secciones posteriores remiten a esta distinción en lugar de redefinirla.

